% 07_conclusion.tex  --  Conclusion

\section{Conclusion}
\label{sec:conclusion}

We introduced the Split-Landscape Discrepancy, a no-candidate-retraining metric
that compares full-data and surrogate GBDT split-gain landscapes in one
histogram pass, and used it to study 17 condensation methods across 26 tabular
datasets. Method-level split-regret correlates with downstream AUROC at
Spearman $\rho=-0.87$, showing that the landscape comparison is a reliable
screening and family-ranking tool before any candidate is retrained.

The central finding overturns a natural intuition: a selector with a provable
$(1-1/e)$ coverage guarantee achieves the lowest downstream AUROC in the
study. Breadth of split-threshold coverage is not a proxy for predictive
fidelity, because downstream performance depends on preserving the
gain-weighted structure of the dominant splits, not on counting reproduced
thresholds. The reference-chain
decomposition (\cref{eq:decomp}) operationalises this distinction: coverage
failure is a structure failure, not a leaf failure, and each failure mode points
to a distinct debugging target, making the decomposition a practical debugging
tool.

The diagnostic reaches its limits at regression tasks, where structure mismatch
dominates and small condensed sets cannot overcome it; at HPO transfer, where
configuration rankings change even when final AUROC is similar; and at very
large scale, where dominant split margins saturate and the score becomes a
rejection filter rather than a fine-grained ranker. Within those boundaries,
the workflow is simple: compute split-regret to screen out poor candidates,
prefer low-regret methods when margins are not saturated, and validate
hyperparameter and regression decisions on full data. We release the benchmark
and implementation to support future evaluations of new condensation methods.
